% Section 9.2, from lectures/L09/appendix/b_freestanding.md.

\section{The Freestanding Target and the Toolchain}
\label{c9:sec:freestanding}\label{c9:app:b}
The host build (g++) links against a full hosted C++ implementation: a heap, exceptions, RTTI,
iostreams, the whole standard library. The ATmega328P has none of that. Its build is
\textbf{freestanding}: avr-gcc, no operating system, and only the thin runtime avr-libc provides.
This section is about what that removes, why the driver survives it unchanged, and how you build and
flash.

\subsection{What freestanding removes, and what replaces it}
\label{c9:sec:removes}
A freestanding avr-gcc build is compiled with
\code{-std=c++14 -mmcu=atmega328p -Os -DF_CPU=16000000UL -fno-exceptions -fno-rtti -fno-threadsafe-statics},
plus \code{-ffunction-sections -fdata-sections} and a linker \code{-Wl,--gc-sections} so that the
linker can drop what virtual dispatch pulls in on a 32\,KB part. Everything the driver uses is C++11,
so the standard here is a floor rather than a requirement; \code{-std=c++14} is chosen because every
avr-gcc in circulation accepts it, while \code{-std=c++17} is not even a recognized option on the
older toolchains, including the one Microchip Studio ships. Three things go away with the rest of
the flags. \textbf{The heap} goes first: avr-libc does provide \code{malloc} and \code{free}, but
there is no C++ \code{operator new} or \code{operator delete}, so dynamic allocation is off the
table in practice and you use static and automatic storage. \textbf{Exceptions and RTTI} go next,
which means no \code{throw}, no \code{dynamic_cast} and no \code{typeid}, and errors become return
values, which is exactly what the driver already does when \code{write} returns \code{false} on a
full FIFO. Finally \textbf{iostreams and most of the standard library} go, so there is no
\code{std::cout}, and logging goes out a real UART instead, as described below.

The payoff of the AVR-portable style of \chapref{6} lands here. Because the register map,
\code{Interface}, \code{Uart}, \code{Stub} and the blocking helpers were all written with
\code{<stdint.h>} and bare types, no \code{std::}, nested namespace blocks, and none of the newer C++
attributes the AVR toolchain rejects, they cross-compile to the target \textbf{untouched}. Nothing
above the seam is rewritten for the AVR; only the transport below it is new. That is the whole point
of having built portable in the first place.

\subsection{\code{env.cpp}, the runtime the compiler still expects}
\label{c9:sec:env}
avr-libc ships no C++ runtime, yet the compiler still emits references to a few runtime symbols for
ordinary C++ constructs, and the link fails without them. The provided \file{avr/env.cpp}
hand-defines three groups of them. \code{operator delete} is needed because a class with a virtual
destructor names the deleting \code{operator delete} in its vtable, even if you never delete through
a base pointer. \code{__cxa_pure_virtual()} is named by every abstract class's vtable, and is called
only if a pure virtual is somehow invoked. And \code{__cxa_guard_acquire()},
\code{__cxa_guard_release()} and \code{__cxa_guard_abort()} form the guard around a function-local
\code{static}'s one-time initialization.

This build passes \code{-fno-threadsafe-statics}, so avr-gcc never emits the \code{__cxa_guard_*}
calls at all; \file{env.cpp} defines them anyway, so that the link succeeds with or without the flag
and on any toolchain that does emit them. \file{env.cpp} lives in \file{avr/}, not \file{source/},
on purpose: the host build already has these from libstdc++, so keeping the file out of the host link
avoids replacing the host's thread-safe static guards with the single-threaded versions. The details
are in \file{fw/README.md}.

\subsection{Logging over the AVR's own UART}
\label{c9:sec:logging}
With no iostreams, debug logging goes out the ATmega328P's \textbf{own hardware USART} (its
\code{TXD} and \code{RXD} pins, the ``serial'' a Nano exposes over USB). Do not confuse it with the
UART peripheral this whole book builds: that one lives on the FPGA and the ATmega drives it over
SPI; this one is the ATmega's built-in serial port, used here purely as a debug console to a PC
terminal.

Using it is a small, poll-based routine: set the baud divider \code{UBRR0} from \code{F_CPU} and the
target baud, enable the transmitter in \code{UCSR0B}, then for each byte wait for the data register
to be empty (\code{UCSR0A & (1U << UDRE0)}) and write it to \code{UDR0}. This is why the logging
path, unlike the SPI transport, needs \code{F_CPU}: the baud divider is computed from the clock.

\subsection{Building and flashing the MCU}
\label{c9:sec:flashing}
Two steps turn the source into something on the chip. First you \textbf{compile and link} with
avr-gcc for \code{-mmcu=atmega328p}, then \code{avr-objcopy} the ELF into an Intel HEX image. Then
you \textbf{flash} it with avrdude over the Nano's USB bootloader. Both are wrapped by the avr-gcc
Makefile in \file{avr/}:

\begin{shell}
make -C fw/avr flash
\end{shell}

\noindent \code{F_CPU} must match the board's actual clock, 16\,MHz on a standard Nano, set by its
crystal, with the fuses selecting that external oscillator as the clock source, because it drives
every computed timing, including the debug UART's baud. The \file{avr/} build is cross-compiled and
flashed only; it is excluded from host CI, while the transport's \emph{logic} is still checked on
the host, as \secref{c9:sec:host_testing} describes.

\subsection{The other toolchain: getting the peripheral onto the FPGA}
\label{c9:sec:quartus}
The FPGA half has the same journey, and this is where it happens, because the two are the same idea
twice: source becomes an image, and an image gets programmed onto a part.

Everything up to now has been GHDL, which only ever simulated your VHDL. Simulation says the design
is \emph{correct}; it says nothing about whether it \textbf{fits} on the device or \textbf{meets
timing} at 50\,MHz. Only synthesis answers those, and the tool for the DE0-CV is \textbf{Quartus
Prime Lite}.

\code{uart_top} is not the top level on the board. It has no notion of which physical pin
\code{sclk} or \code{tx} is, so it is wrapped by \textbf{\file{uart_board.vhd}}, the provided
Quartus top level that maps its ports onto DE0-CV package pins and feeds it the board's 50\,MHz
clock. That wrapper is board I/O rather than peripheral logic, which is why it lives with the
Quartus project instead of in \file{hw/}.

The flow is: open the project, add your \file{hw/*.vhd} alongside the wrapper, \textbf{compile},
then read the two numbers that matter. \textbf{Fit} tells you the design is small enough, which for
this peripheral it comfortably is. \textbf{Timing} tells you every path from flip-flop to flip-flop
settles inside one 20\,ns clock period; a design that simulates perfectly can still fail here, and
that failure is invisible to every testbench you have run. Then \textbf{program} the board over
USB-Blaster.

What the board can show with nothing else attached is less than it may seem. \code{uart_top} has no
path of its own from \code{rx} to \code{tx}: it transmits only what an SPI master writes to
\code{TX_DATA}, and it holds what it receives in the RX FIFO until a master reads it, so with no
Nano on the SPI pins \code{BAUD_DIV} reads zero and \code{tx} simply idles high. The loopback of
\code{uart_top_tb} worked because the testbench was that master. What can be checked alone is the
path around the peripheral: wire the wrapper's \code{rx} pin straight back out to its \code{tx} pin,
bypassing \code{uart_top}, re-synthesize, and a terminal on the USB-serial adapter echoes what you
type. That proves the programming, the pin assignment, the adapter and its logic levels
(not the terminal's baud rate: the adapter sends and receives at the same setting, so a loopback agrees with itself at any rate), and it is the first rung of the bring-up ladder of \chapref{10}; the peripheral
itself is first exercised there, once the Nano writes \code{BAUD_DIV} and \code{TX_DATA} over SPI.

\subsection{What's ahead}
The exercises in \secref{c9:sec:exercises} implement \code{AvrSpi}, write the freestanding bring-up
\code{main}, add UART logging, and measure a register round trip against the 1\,MHz \code{SCK}
budget. \chapref{10} then puts both halves on the bench.
